%%%%%%%%%%%%%%%%
% RESUMEN
%%%%%%%%%%%%%%%%

\thispagestyle{empty}
\begin{center}
    \Large{Ficha Resumen \\ Anteproyecto de Trabajo de Grado}
\end{center}

\textbf{Posible Título:} Desarrollo de una herramienta CLI para la generación automatizada de proyectos backend e integración de módulos reutilizables
\begin{enumerate}
    \item Área de trabajo: Automatización del Desarrollo y Reutilización de Software
    \item Tipo de proyecto: Aplicado
    \item Estudiante: Jorge Andres Facundo Solarte
    \item Correo electrónico: jfacundo@javerianacali.edu.co
    \item Dirección y teléfono: Calle 2\#19-250, Jamundí, Valle del Cauca
    \item Director: Juan Carlos Martinez Arias
    \item Vinculación del director: Javeriana Cali, Planta
    \item Correo electrónico del director: % PENDIENTE: correo de Juan Carlos Martinez Arias
    \item ODS que aplica al proyecto (Agenda 2030): ODS 9 - Industria, Innovación e Infraestructura
    \item Palabras clave: Reutilización de software, Scaffolding de proyectos, Arquitectura hexagonal, Domain-Driven Design, BIAN
    %\item Fecha de inicio:
    \item Duración estimada: cinco meses
    \item Resumen:  El proyecto aborda la necesidad de reducir el esfuerzo repetitivo y las diferencias estructurales entre proyectos que enfrentan los equipos de desarrollo al iniciar nuevos proyectos de software e integrar funcionalidades compartidas en entornos bancarios. Se propone el diseño, implementación y validación de una herramienta CLI para entornos Node.js, orientado al desarrollo de aplicaciones backend en el sector financiero, que automatice el arranque de proyectos e integre capacidades reutilizables de forma desacoplada, alineado con estándares bancarios como BIAN e ISO 20022.
\end{enumerate}